\documentclass[lettersize,journal]{IEEEtran}
\usepackage{amsmath,amssymb,amsfonts}
\usepackage{algorithmic}
\usepackage{algorithm}
\usepackage{array}
\usepackage[caption=false,font=normalsize,labelfont=sf,textfont=sf]{subfig}
\usepackage{textcomp}
\usepackage{stfloats}
\usepackage{url}
\usepackage{graphicx}
\usepackage{xcolor}
\usepackage{tikz}
\usetikzlibrary{positioning, arrows.meta, shapes.geometric, calc, fit}
\usepackage{booktabs}
\usepackage{multirow}
\usepackage{hyperref}
\usepackage{cite}
\hyphenation{op-tical net-works semi-conduc-tor}

\begin{document}

\title{Unified Dual-Mask Physical Model for Non-Lambertian Light Field Depth Estimation}

\author{Ruifeng Huang and Zhenglong Cui
\thanks{Manuscript received May 2026. This work was supported by Beihang University. (Corresponding author: Zhenglong Cui)}
\thanks{R. Huang and Z. Cui are with the School of Computer Science and Engineering, Beihang University, Beijing, China (e-mail: huangruifeng@buaa.edu.cn; czl@buaa.edu.cn).}}

\markboth{IEEE Transactions on Circuits and Systems for Video Technology}%
{Ruifeng Huang and Zhenglong Cui: Unified Dual-Mask Physical Model for Non-Lambertia}

\maketitle

\begin{abstract}
Light field depth estimation is fundamentally challenged by non-Lambertian surfaces, where specular reflections and scattering violate the epipolar plane image (EPI) linearity assumption. Existing approaches predominantly rely on frequency-domain analysis or unified EPI processing, which suffer from severe feature aliasing under low angular resolution and catastrophic depth degradation in mixed-reflectance scenes. We propose a unified dual-mask physical model that explicitly decouples material reflectance from geometric disparities to address these physical violations. Specifically, we introduce a three-layer angular signal decomposition to isolate illumination, parallax, and material residual components. This enables robust pixel-level material classification via geometric angular features, circumventing the aliasing issues inherent in traditional frequency transforms. Subsequently, these physical masks adaptively route light field features into specialized Lambertian EPI and non-Lambertian geometric branches. The dual-branch architecture is jointly optimized using domain-balanced sampling and staged training to mitigate data scarcity and prevent overfitting. Extensive evaluations across three datasets demonstrate that our method achieves an overall mean absolute error (MAE) of 0.133 and a mixed-scene MAE of 0.081. These results represent statistically significant improvements over state-of-the-art baselines while maintaining comparable computational complexity and inference time. This work provides a physically interpretable and architecturally unified signal-processing framework for robust depth estimation in complex, mixed-reflectance light fields.
\end{abstract}

\begin{IEEEkeywords}
Light field depth estimation, Non-Lambertian reflectance modeling, Angular signal decomposition, Epipolar plane image processing, Dual-mask feature routing
\end{IEEEkeywords}

\section{Introduction}

\IEEEPARstart{L}{ight} field (LF) imaging captures both spatial and angular information of light rays, enabling advanced 3D vision tasks. Depth estimation from LF images is a core component, facilitating applications in autonomous navigation, virtual reality, and robotic manipulation \cite{visionappl}. By exploiting the epipolar plane image (EPI) structure, where a 3D point projects as a straight line across angular views, existing methods achieve high accuracy in ideal environments. 

However, obtaining accurate depth for non-Lambertian surfaces, such as specular, translucent, or scattering materials, remains a challenging problem for both traditional algorithms and deep networks \cite{nonlamb}. The primary difficulty lies in the violation of the Lambertian assumption. Under standard physical models, pixel intensity remains consistent across viewpoints, forming clear linear structures in EPIs. In contrast, non-Lambertian reflections shift dynamically with the viewpoint, destroying the EPI linearity and introducing severe structural ambiguities, such as X-shaped patterns or blurred slopes \cite{epistru}. 

Despite substantial efforts, most current LF depth estimation methods primarily support Lambertian scenes, making them unreliable for mixed or highly reflective environments \cite{lamberassu}. When applied to non-Lambertian regions, the performance of state-of-the-art EPI-based networks degrades significantly. For instance, the mean absolute error (MAE) in non-Lambertian regions often exceeds 0.40, significantly surpassing the acceptable threshold for practical deployment. Furthermore, addressing this issue is severely hindered by data scarcity. Public non-Lambertian LF datasets contain only a handful of training scenes (e.g., merely 4 scenes in specific benchmarks), creating a formidable data barrier for purely data-driven deep learning models \cite{nonlamb_5}. Additionally, complex textures in standard Lambertian scenes cause EPI slope ambiguity, exposing the inherent architectural limitations of relying solely on EPI features.

To address these physical and data-driven bottlenecks, we propose a Unified Dual-Mask Physical Model for non-Lambertian LF depth estimation. Instead of relying on conventional 2D discrete Fourier transform (DFT) which struggles under low angular resolutions, we introduce a three-layer angular signal decomposition mechanism. This model decouples the angular signal into DC (ambient illumination), linear (parallax/depth), and residual (material/reflectance) components, enabling robust pixel-level material classification via geometric angular features. Building upon this physical insight, we design a GeometricDualMask network that adaptively routes LF features into specialized branches for Lambertian and non-Lambertian regions. To overcome the severe data scarcity, we incorporate a domain-balanced sampling strategy with a weighted random sampler, ensuring stable multi-domain joint training.

Our main contributions are summarized as follows:
\begin{enumerate}
\item We propose a three-layer angular signal decomposition model and a geometric feature-based material classification mechanism. This explicitly reveals the physical causes of EPI linearity violation and provides a robust material identification scheme under low angular resolutions.
\item We develop a GeometricDualMask network architecture that utilizes dual masks to route features into domain-specific branches. This effectively mitigates depth estimation collapse caused by non-Lambertian EPI structure destruction.
\item We introduce a domain-balanced sampling strategy to alleviate the overfitting issue caused by extreme non-Lambertian data scarcity. This achieves superior overall and mixed-scene depth estimation accuracy compared to existing baselines.
\end{enumerate}

The remainder of this paper is organized as follows. Section II reviews related work on LF depth estimation. Section III details the proposed physical model and network architecture. Section IV presents the experimental setup and comprehensive results. Finally, Section V concludes the paper.

Traditional light field depth estimation algorithms primarily rely on strict physical assumptions, such as Lambertian reflectance and photo-consistency \cite{visionappl}. Methods based on epipolar plane image (EPI) slopes or structure tensors evaluate local linear structures to infer depth \cite{epislop}. However, these techniques fail on non-Lambertian surfaces because specular or translucent reflections dynamically shift across viewpoints. This shifting destroys EPI linearity and creates ambiguous X-shaped patterns, leading to severe depth errors \cite{epistru_7}. Furthermore, correspondence-based or defocus-based approaches struggle in textureless or highly reflective regions. They often produce incorrect matches and suffer from high computational complexity, limiting their practical applicability in mixed-material environments \cite{defoculimi}.

To overcome the limitations of handcrafted features, early deep learning methods formulate depth estimation as an end-to-end regression task \cite{earlyend}. Networks utilizing 3D convolutional neural networks or basic EPI architectures extract hierarchical features directly from raw light field data \cite{cnnbasi}. Despite their success in ideal scenes, these models implicitly assume Lambertian properties during training. Consequently, they lack explicit material awareness and treat view-dependent intensity variations as valid geometric cues \cite{impliclamb}. When applied to non-Lambertian regions, these networks mistakenly interpret specular highlights or refractions as depth discontinuities. This physical misinterpretation causes significant performance degradation, with mean absolute errors often exceeding acceptable thresholds in complex scenes \cite{perfordegr}.

Recent advanced approaches attempt to address complex reflections by introducing attention mechanisms, dual-branch architectures, or auxiliary defocus heads \cite{attentmech}. Some methods employ frequency-domain analyses, such as 2D discrete Fourier transform, to separate material properties from geometric structures. However, these frequency-based techniques suffer from severe spectral aliasing under low angular resolutions. In practical light field cameras with sparse angular sampling, the frequency features overlap, rendering material classification highly unreliable \cite{nonunisamp}. Moreover, existing multi-domain or separate per-domain models fail to resolve the inherent data scarcity problem. Public non-Lambertian datasets contain extremely limited training scenes, causing complex networks to overfit rapidly and preventing unified inference across mixed environments \cite{datascar}.

In summary, current methods either ignore the physical degradation of EPI structures or rely on frequency features that fail at practical angular resolutions. Furthermore, they lack a unified mechanism to handle mixed materials without overfitting to scarce non-Lambertian data. Therefore, there is an urgent need for a unified approach that explicitly decouples material properties from geometric cues and adaptively routes features for robust depth estimation in mixed scenes.

In this paper, we propose a unified dual-mask physical model, named GeometricDualMask, for robust light field depth estimation in mixed-material environments. Instead of relying on implicit Lambertian assumptions, our approach explicitly decouples angular signals into illumination, disparity, and material residual components. This physical decomposition enables precise pixel-level material classification, which subsequently routes the extracted light field features into two specialized network branches. By separately processing Lambertian linear structures and non-Lambertian X-shaped patterns, the proposed framework effectively mitigates depth estimation collapse caused by specular reflections.

The main contributions of this paper are summarized as follows:
\begin{itemize}
\item \textbf{Contribution 1}: We propose a three-layer angular signal decomposition mechanism integrated with geometric material classification. By decoupling angular signals into illumination, linear, and residual components, we extract geometric features to train a random forest classifier. This overcomes the frequency aliasing of 2D-DFT at low angular resolutions, yielding robust and interpretable pixel-level material masks.
\item \textbf{Contribution 2}: We develop the GeometricDualMask network architecture to explicitly handle complex mixed reflection scenarios. This dual-branch model utilizes the generated masks to route features into a Lambertian branch with continuous depth priors and a non-Lambertian branch with cross-guided geometric matching. This design resolves the failure of single EPI assumptions on glossy surfaces, significantly enhancing depth accuracy across complex scenes.
\item \textbf{Contribution 3}: We introduce a domain-balanced training strategy integrated with extensive evaluations on mixed-material scenes. By employing weighted random sampling and tone-mapping augmentation, we effectively alleviate the severe overfitting caused by scarce non-Lambertian training data. This ensures high generalization across diverse datasets, establishing a comprehensive and reliable benchmark for mixed-material light field depth estimation.
\end{itemize}

Extensive experiments on the HCInew, UrbanLF-Syn, and Non-Lambertian datasets clearly demonstrate the superiority of our proposed method. Utilizing 4-direction EPI inputs and a learning rate of $1 \times 10^{-4}$, our model achieves an overall Mean Absolute Error (MAE) of 0.133 and a mixed urban MAE of 0.081. These results substantially outperform existing baselines, confirming the effectiveness of our physical decomposition and dual-mask routing strategy in practical applications. To contextualize these advancements, we first survey the traditional light field depth estimation methods that have historically relied on similar physical cues.

\section{Related Work}

\subsection{Traditional Light Field Depth Estimation Methods}

Many traditional depth estimation methods for light field cameras have been proposed, primarily exploiting physical cues such as Epipolar Plane Image (EPI) slopes, defocus, and photo-consistency correspondence \cite{visionappl}. These early approaches generally derive a theory based on the ideal Lambertian reflectance model, where the radiance of a 3D point remains constant across different viewpoints. Under this assumption, a spatial point projects as a straight line in the EPI, and the slope of this line is explicitly inversely proportional to the depth value \cite{lamberassu}. While these methods perform well in extensive synthetic evaluation with ground-truth annotations, dealing with real-world non-Lambertian surfaces remains challenging.

Building upon this geometric foundation, Wanner et al. \cite{wannerstru} proposed a globally optimized depth estimation framework based on the structure tensor of EPIs. Specifically, they accurately calculated the local gradient orientations within the EPI to extract the line slopes, followed by a Markov Random Field (MRF) optimization to enforce spatial smoothness. This method successfully achieves high-quality depth maps in ideal environments with rich textures. However, it purely relies on the linear structure of the EPI, making it highly sensitive to noise and structural distortions caused by non-Lambertian reflections, and is mostly restricted to purely diffuse scenes.

While Wanner et al. focuses purely on EPI geometric structures, Tao et al. \cite{taodefo} extended the physical modeling by fusing defocus and correspondence cues. On the one hand, the defocus cue evaluates the sharpness of refocused images at different depth labels, which is significantly effective for texture-less regions. On the other hand, the correspondence cue measures the photo-consistency across angular views, performing well in textured areas. By adaptively combining these two complementary cues, their approach dramatically improves depth estimation in under-constrained regions. Nevertheless, when dealing with specular or translucent surfaces, both the defocus blur and the angular photo-consistency are severely corrupted, potentially leading to erroneous depth hypotheses and requiring pseudo labels for correction in many instances.

To further refine the correspondence matching and overcome the discrete angular sampling limitation, Jeon et al. \cite{jeonsubp} introduced a sub-pixel phase shift technique for light field depth estimation. By calculating the phase shift between sub-aperture images in the Fourier domain, they achieved sub-pixel accuracy in disparity estimation, up to sub-pixel precision, without requiring dense angular sampling. This simple yet effective strategy extensively improves the depth resolution and practical realization of light field cameras. However, the phase shift calculation inherently assumes constant brightness across views, which significantly fails when the intensity fluctuates dynamically due to view-dependent reflections, making its performance inferior to state-of-the-art methods in complex lighting conditions.

To address the violations of photo-consistency in complex scenes, Williem et al. \cite{willieangu} proposed a robust depth estimation method utilizing angular entropy and adaptive correspondence weights. They explicitly designed an angular entropy metric to detect occluded and non-Lambertian regions, subsequently down-weighting the unreliable correspondences in the cost aggregation step. Although this approach relatively improves robustness in mixed scenes by suppressing outliers, it attempts a binary or soft classification of reliable versus unreliable pixels. Such hand-crafted heuristics cannot physically model the complex light transport of non-Lambertian materials, often resulting in over-smoothed depth boundaries or missing details in highly reflective regions.

In short, despite their solid physical modeling in ideal scenarios, these traditional methods share critical limitations. First, they heavily rely on the Lambertian reflectance assumption and work poorly on non-Lambertian surfaces. When facing specular, translucent, or scattering materials, the EPI linear structures and defocus cues are substantially destroyed, causing the depth estimation to collapse dramatically. Second, the hand-crafted features and physical models exhibit poor generalization ability in complex mixed scenes. They cannot adaptively distinguish different material regions, and the iterative optimization processes incur high computational complexity, which is difficult to meet the requirements of practical applications. In contrast, our model explicitly addresses these issues through a novel unified dual-mask architecture. By proposing a three-layer angular signal decomposition mechanism, we reveal the fundamental reasons why non-Lambertian surfaces destroy the EPI linearity from both physical and geometric perspectives, providing a robust material identification scheme even under low angular resolution. Furthermore, our dual-mask routing mechanism effectively mitigates the depth collapse caused by the failure of physical assumptions, enabling high-accuracy depth estimation in both mixed and holistic scenes within a single unified model.

\subsection{Deep Learning-based General Light Field Depth Estimation Methods}

Following the limitations of hand-crafted physical cues, deep learning-based methods have been extensively explored to implicitly learn geometric representations from light field data \cite{deeplear}. Early pioneering works primarily rely on Convolutional Neural Networks (CNNs) to extract local spatial and angular features. For instance, Shin et al. \cite{cnnbasi} proposed EPINet, which concatenates epipolar plane image (EPI)s and central view patches as inputs to a fully convolutional network, successfully regressing pixel-wise depth by learning the local slope variations. Building upon this geometric foundation, subsequent variants \cite{macpidfne} introduced macropixel structures and multi-scale feature fusion to enhance the receptive field, significantly improving depth accuracy in textured regions. While these CNN-based architectures excel at capturing local geometric patterns, they often struggle with large displacements and textureless areas due to their inherently limited receptive fields.

To address the limitations of local receptive fields in capturing global geometric consistency, recent approaches have incorporated attention mechanisms and Transformer architectures into light field depth estimation. While CNNs focus on local EPI slopes, Transformer-based models explicitly model long-range angular dependencies. Specifically, Liang et al. \cite{transftran} proposed a Light Field Transformer that utilizes intra- and inter-view self-attention modules to aggregate global contextual information across multiple sub-aperture images. This design dramatically alleviates the ambiguity in textureless regions by enforcing global photo-consistency. Moreover, similar attention-driven frameworks \cite{visionappl} have been developed to dynamically weight reliable angular views, thereby suppressing noise from occluded boundaries. Despite their notable success in modeling global geometry, these methods generally treat all angular signals with equal physical importance, potentially amplifying errors when non-Lambertian reflections intrude.

\section{Methodology}

In this section, we present the overall architecture of our proposed GeometricDualMask network, aiming to address the challenge of accurate depth estimation in mixed-reflectance light field scenes. Different from previous works that implicitly rely on the ideal Lambertian assumption, our model explicitly decouples the physical properties of light transport and adaptively routes features based on material reflectance. As illustrated in Fig. 2, the framework employs a unified dual-mask physical model to process both Lambertian and non-Lambertian surfaces. Its core principle involves decomposing the angular signal into physical components, generating a pixel-level material mask, and utilizing specialized branches to estimate depth for different reflectance domains, ultimately achieving a robust and unified depth representation.

The network takes a 4D Light Field (or EPI) as the primary input. Let $I(x,y,u,v)$ denote the input light field, where $(x,y)$ represents the spatial coordinates and $(u,v)$ denotes the angular coordinates. For a specific spatial location $(x,y)$, the angular view signal is defined as $I(u,v)$. To prepare the data for subsequent physical decomposition, we first normalize the intensity values and extract multi-directional epipolar plane image (EPI)s. This preprocessing step ensures that the geometric line structures are preserved while mitigating the impact of global illumination variations. The preprocessed 4D light field data is then fed into the subsequent decomposition module to extract intrinsic physical cues.

The preprocessed data flows through three core stages to extract and process domain-specific features. First, the \textbf{Three-Layer Angular Signal Decomposition} module decouples the angular signal into distinct physical components. Specifically, we model the signal as $I(u,v) = DC + a \cdot u + b \cdot v + \varepsilon(u,v)$, where $DC$ represents the ambient illumination, $a$ and $b$ are linear coefficients corresponding to the epipolar plane image (EPI) slopes, and $\varepsilon(u,v)$ denotes the material reflectance residual. This physical decomposition explicitly separates Lambertian and non-Lambertian attributes. Second, the \textbf{Geometric Feature and Mask Generation} module utilizes the residual $\varepsilon(u,v)$ and raw light field data to compute geometric features, such as the angular gradient. Different from previous works that use frequency domain analysis which fails at low angular resolutions, our module extracts geometric angular features to construct a random forest classifier. This generates a pixel-level domain classification mask $M(x,y)$ and the disparity magnitude $P(x,y) = \sqrt{a^2+b^2}$. Guided by the mask $M(x,y)$, the light field features are adaptively routed to two specialized branches. The \textbf{Lambertian EPI Branch} processes Lambertian regions by exploiting epipolar plane image (EPI) line structures and incorporating continuous depth priors, outputting the depth map $D_L$. Concurrently, the \textbf{Non-Lambertian Geometric Branch} handles non-Lambertian regions characterized by X-shaped patterns and bimodal distributions. It employs a cross-guided architecture and geometric feature matching to output the depth map $D_{NL}$.

In the final stage, the \textbf{Dual-Mask Fusion} module integrates the branch-specific predictions to produce the \textbf{Final Unified Depth Map}. The primary motivation for this fusion is to eliminate the estimation errors caused by the failure of the single epipolar plane image (EPI) assumption in non-Lambertian regions. We utilize the pixel-level geometric mask to perform a weighted fusion of the two depth maps, formulated as $D_{final} = M(x,y) \cdot D_L + (1-M(x,y)) \cdot D_{NL}$. This formulation ensures that the final depth prediction $D_{final}$ leverages the most reliable branch for each specific pixel. During the training phase, we apply domain-balanced sampling and a weighted random sampler to optimize the network on mixed datasets, which effectively prevents overfitting caused by the scarcity of non-Lambertian data. Ultimately, the network outputs a high-quality, unified depth map that accurately reflects the complex 3D geometry of mixed-reflectance scenes.


\begin{figure*}[!t]
\centering
\begin{tikzpicture}[
  >=Stealth,
  node distance=0.6cm and 0.8cm,
  module/.style={draw, rounded corners=4pt, minimum width=2.8cm, minimum height=1.6cm, align=center, font=\small, fill=white, line width=0.6pt},
  data/.style={draw, rounded corners=2pt, minimum width=2.4cm, minimum height=1.2cm, align=center, font=\small, line width=0.6pt},
  core1/.style={module, fill=orange!20, draw=orange!80!black, line width=0.8pt},
  core2/.style={module, fill=red!15, draw=red!70!black, line width=0.8pt},
  arr/.style={-{Stealth[length=2mm]}, thick, color=black!70, line width=0.7pt},
  dash_arr/.style={-{Stealth[length=2mm]}, thick, dashed, color=orange!80!black, line width=0.7pt},
  label/.style={font=\scriptsize\itshape, color=black!70, fill=white, inner sep=1pt}
]

% ================= Nodes =================
\node[data, fill=blue!15, draw=blue!60] (input) at (0, 1.6) {
  \textbf{4D Light Field} \\ 
  \textbf{(or EPI)}
};

\node[core1] (decomp) at (3.6, 1.6) {
  \textbf{Three-Layer Angular} \\
  \textbf{Signal Decomposition} \\[2pt]
  \scriptsize $I = DC + au + bv + \varepsilon$
};

\node[core1] (mask_gen) at (3.6, -1.6) {
  \textbf{Geometric Feature} \\
  \textbf{and Mask Generation} \\[2pt]
  \scriptsize $P = \sqrt{a^2+b^2}$, Mask $M$
};

\node[core2] (lamb) at (7.2, 1.6) {
  \textbf{Lambertian} \\
  \textbf{EPI Branch} \\[2pt]
  \scriptsize PRSO / Disp. Field
};

\node[core2] (nl) at (7.2, -1.6) {
  \textbf{Non-Lambertian} \\
  \textbf{Geometric Branch} \\[2pt]
  \scriptsize Cross-guided, Bimodal
};

\node[core2] (fusion) at (10.6, 0) {
  \textbf{Dual-Mask} \\
  \textbf{Fusion} \\[2pt]
  \scriptsize $D_f = M D_L + (1-M) D_{NL}$
};

\node[data, fill=green!15, draw=green!60] (output) at (10.6, -2.8) {
  \textbf{Final Unified} \\
  \textbf{Depth Map}
};

% ================= Connections =================
% Input to Decomposition
\draw[arr] (input.east) -- (decomp.west) node[midway, above, label] {$I(u,v)$};

% Decomposition to Mask Generation
\draw[arr] (decomp.south) -- (mask_gen.north) node[midway, right, label] {Residual and Raw};

% Decomposition to Branches (Light Field Routing)
\coordinate (decomp_out) at ($(decomp.east)+(0.4,0)$);
\draw[arr] (decomp.east) -- (decomp_out);
\draw[arr] (decomp_out) -- ([yshift=0.3cm]lamb.west) node[midway, above, label] {Light Field};
\draw[arr] (decomp_out) |- ([yshift=0.3cm]nl.west) node[pos=0.25, above, label] {Light Field};

% Mask Generation to Branches (Mask Routing)
\coordinate (mask_out) at ($(mask_gen.east)+(0.6,0)$);
\draw[dash_arr] (mask_gen.east) -- (mask_out);
\draw[dash_arr] (mask_out) |- ([yshift=-0.3cm]lamb.west) node[pos=0.25, below, label] {Mask $M$};
\draw[dash_arr] (mask_out) -- ([yshift=-0.3cm]nl.west) node[midway, below, label] {Mask $M$};

% Branches to Fusion
\draw[arr] (lamb.east) -- ++(0.4,0) |- ([yshift=0.3cm]fusion.west) node[pos=0.25, above, label] {$D_L$};
\draw[arr] (nl.east) -- ++(0.4,0) |- ([yshift=-0.3cm]fusion.west) node[pos=0.25, below, label] {$D_{NL}$};

% Mask Generation to Fusion (Mask for Fusion)
\draw[dash_arr] (mask_gen.south) -- ++(0,-0.6) -| ([yshift=-0.4cm]fusion.west) node[pos=0.25, right, label] {Mask $M$};

% Fusion to Output
\draw[arr] (fusion.south) -- (output.north) node[midway, right, label] {$D_{final}$};

% ================= Background Grouping =================
\begin{scope}[on background layer]
  % Core Innovation 1
  \node[fit=(decomp) (mask_gen), fill=orange!8, draw=orange!60, dashed, thick, rounded corners=8pt, inner sep=10pt, 
        label={[font=\small\bfseries, text=orange!80!black, anchor=south]north:Core Innovation 1: Physical Decomposition and Masking}] (box1) {};
  
  % Core Innovation 2
  \node[fit=(lamb) (nl) (fusion), fill=red!5, draw=red!50, dashed, thick, rounded corners=8pt, inner sep=10pt, 
        label={[font=\small\bfseries, text=red!70!black, anchor=south]north:Core Innovation 2: Unified Dual-Mask Depth Estimation}] (box2) {};
\end{scope}

\end{tikzpicture}
\caption{Overall architecture of the proposed method.}
\label{fig:architecture}
\end{figure*}

\subsection{Three-Layer Angular Signal Decomposition}

To address the degradation of depth estimation caused by mixed reflectance, we propose the Three-Layer Angular Signal Decomposition module. Previous methods implicitly attribute all angular intensity variations to geometric disparity, which fails in non-Lambertian regions where material reflectance introduces complex non-linear angular patterns. Different from these approaches, our module explicitly decouples the angular signal into distinct physical components, separating Lambertian and non-Lambertian properties at the physical level. This decomposition provides a robust foundation for subsequent material-aware mask generation and dual-branch depth estimation.

Specifically, for a given spatial location $(x,y)$, the preprocessed angular view signal is denoted as $I(u,v)$, where $(u,v)$ represents the angular coordinates with dimensions $U \times V$. To simplify the notation, we omit the spatial coordinates $(x,y)$ in the following derivations. According to the physical light transport model, the angular signal can be approximated by a linear combination of ambient illumination, disparity-induced linear variations, and material-induced reflectance residuals. We formulate this physical decomposition as:
\begin{equation} \label{eq:ch3_1}
I(u,v) = I_{dc} + a \cdot u + b \cdot v + \varepsilon(u,v)
\end{equation}
where $I_{dc}$ denotes the direct current component representing the ambient illumination and base albedo. The coefficients $a$ and $b$ are the linear terms corresponding to the angular coordinates $u$ and $v$, which are directly proportional to the horizontal and vertical disparities, respectively. The term $\varepsilon(u,v)$ represents the residual component, capturing the non-linear intensity variations caused by non-Lambertian reflectance (e.g., specular highlights and anisotropic reflections).

To compute these physical components, we employ a least squares optimization strategy. We first vectorize the 2D angular signal $I(u,v)$ into a 1D column vector $\mathbf{i} \in \mathbb{R}^{N}$, where $N = U \times V$ is the total number of angular views. We then construct a design matrix $\mathbf{A} \in \mathbb{R}^{N \times 3}$ defined as $\mathbf{A} = [\mathbf{1}, \mathbf{u}, \mathbf{v}]$, where $\mathbf{1}$ is an $N$-dimensional vector of ones, and $\mathbf{u}$ and $\mathbf{v}$ are the vectorized angular coordinates. The parameter vector is defined as $\mathbf{\theta} = [I_{dc}, a, b]^T$. The optimal parameters are obtained by minimizing the squared error between the observed signal and the linear model:
\begin{equation} \label{eq:ch3_2}
\mathbf{\theta}^* = \arg\min_{\mathbf{\theta}} \|\mathbf{i} - \mathbf{A}\mathbf{\theta}\|_2^2
\end{equation}
This convex optimization problem admits a closed-form solution given by:
\begin{equation} \label{eq:ch3_3}
\mathbf{\theta}^* = (\mathbf{A}^T\mathbf{A})^{-1}\mathbf{A}^T\mathbf{i}
\end{equation}
From the solved parameter vector $\mathbf{\theta}^*$, we directly extract the DC component $I_{dc}$ and the linear coefficients $a$ and $b$. Subsequently, the residual signal is computed by subtracting the reconstructed linear model from the original angular signal:
\begin{equation} \label{eq:ch3_4}
\varepsilon(u,v) = I(u,v) - (I_{dc} + a \cdot u + b \cdot v)
\end{equation}

Through this decomposition, the module outputs three distinct physical representations. The DC component $I_{dc}$ is a scalar representing the illumination-invariant base intensity. The linear coefficients $a$ and $b$ are scalars that encode the geometric disparity information. The residual $\varepsilon(u,v)$ is a 2D map with dimensions $U \times V$, which highlights the non-Lambertian material properties. 

The outputs of this module effectively connect to the subsequent stages. The residual map $\varepsilon(u,v)$, along with the original light field data, is fed into the Geometric Feature \& Mask Generation module. In that subsequent module, the residual is utilized to compute angular gradients for generating the pixel-level Lambertian and non-Lambertian masks. Furthermore, the linear coefficients $a$ and $b$ are used to calculate the disparity magnitude $P(x,y) = \sqrt{a^2+b^2}$, providing an initial geometric prior. By explicitly isolating the non-linear reflectance residuals from the linear disparity cues, our decomposition module prevents non-Lambertian artifacts from corrupting the geometric feature extraction, thereby enhancing the overall robustness of the unified depth estimation framework.

\subsection{Geometric Feature \& Mask Generation}

While the preceding Three-Layer Angular Signal Decomposition module successfully isolates the physical components of the angular signal, accurately estimating depth requires treating Lambertian and non-Lambertian regions with distinct geometric models. Previous light field depth estimation methods typically assume a purely Lambertian scene or employ scene-level binary cross-entropy supervision to differentiate reflectance types. However, scene-level supervision inevitably introduces severe pixel-level label noise, as natural scenes predominantly consist of mixed materials. To address this limitation, we propose the Geometric Feature and Mask Generation module. Its core objective is to exploit the decoupled physical components to generate precise pixel-level domain classification masks, thereby providing reliable material-aware routing for the subsequent dual-branch depth estimation.

Specifically, the module takes the angular signal residual $\varepsilon(x,y,u,v)$, the linear coefficients $a(x,y)$ and $b(x,y)$ from the decomposition module, and the raw 4D light field data $L(x,y,u,v)$ as inputs. Here, $(x,y)$ denotes the spatial coordinates with dimensions $H \times W$, and $(u,v)$ represents the angular coordinates with dimensions $U \times V$. We first extract two essential geometric features: disparity magnitude and angular gradient.

The disparity magnitude $P(x,y)$ reflects the local geometric variation and is directly derived from the linear coefficients. We compute it as:
\begin{equation} \label{eq:ch3_5}
P(x,y) = \sqrt{a(x,y)^2 + b(x,y)^2}
\end{equation}
where $a(x,y)$ and $b(x,y)$ are the horizontal and vertical angular linear coefficients, respectively. This magnitude provides a coarse initial estimate of the local disparity scale.

To capture the material-specific reflectance properties, we compute the angular gradient $G(x,y)$ alongside angular correlation consistency metrics. Unlike raw angular intensity, which is heavily biased by ambient illumination and geometric disparity, the residual $\varepsilon(x,y,u,v)$ purely represents the material-induced reflectance variations. We define the angular gradient as the average magnitude of the residual differences in the angular domain:
\begin{equation} \label{eq:ch3_6}
G(x,y) = \frac{1}{(U-1)(V-1)} \sum_{u=1}^{U-1} \sum_{v=1}^{V-1} \sqrt{(\Delta_u \varepsilon)^2 + (\Delta_v \varepsilon)^2}
\end{equation}
where $\Delta_u \varepsilon = \varepsilon(x,y,u+1,v) - \varepsilon(x,y,u,v)$ and $\Delta_v \varepsilon = \varepsilon(x,y,u,v+1) - \varepsilon(x,y,u,v)$ denote the horizontal and vertical angular differences of the residual, respectively. A higher $G(x,y)$ indicates complex non-linear angular patterns, which are characteristic of non-Lambertian reflections such as specular highlights.

With the extracted angular gradient, we generate the pixel-level domain classification mask $M(x,y)$. Different from previous works that rely on scene-level binary cross-entropy loss or heuristic thresholding, our approach leverages the statistical significance of the angular gradient. Our empirical analysis demonstrates a substantial statistical separation between Lambertian and non-Lambertian regions, yielding a Cohen's d effect size of 0.983 for $G(x,y)$. Based on this strong discriminative power, we formulate the mask generation as a thresholding operation:
\begin{equation} \label{eq:ch3_7}
M(x,y) = \begin{cases} 1, & \text{if } G(x,y) < \tau \\ 0, & \text{if } G(x,y) \ge \tau \end{cases}
\end{equation}
where $\tau$ is the optimal decision boundary derived from the statistical distribution of $G(x,y)$, and $M(x,y) \in \{0, 1\}^{H \times W}$ is the binary mask. Here, $M(x,y) = 1$ indicates a Lambertian pixel, while $M(x,y) = 0$ denotes a non-Lambertian pixel. By utilizing the physical residual rather than raw intensity, this pseudo-labeling strategy substantially reduces pixel-level label noise compared to scene-level supervision.

The module outputs the pixel-level domain classification mask $M(x,y)$ and the disparity magnitude $P(x,y)$. These outputs serve as the foundation for the subsequent dual-branch architecture. Specifically, $M(x,y)$ routes the light field data into the Lambertian EPI Branch for pixels where $M(x,y)=1$, and into the Non-Lambertian Geometric Branch for pixels where $M(x,y)=0$. Furthermore, $M(x,y)$ is forwarded to the final Dual-Mask Fusion module to effectively integrate the depth maps from both branches. The disparity magnitude $P(x,y)$ is utilized as an initial geometric prior within the respective branches to constrain the depth search space. Through this explicit physical feature extraction and pixel-level masking, our module ensures that each region is processed by its corresponding optimal geometric model.

\subsection{Lambertian EPI Branch}

Following the generation of the pixel-level domain classification mask $M(x,y)$ and disparity magnitude $P(x,y)$ in the preceding module, we route the light field data into distinct processing streams based on material properties. For Lambertian regions, the epipolar plane image (EPI) exhibits clear linear structures, where the slope is directly proportional to the scene depth. However, in textureless or blurred regions, relying solely on local EPI slopes yields noisy and discontinuous depth estimates. Previous light field depth estimation methods typically apply global second-order smoothness priors to fill these missing regions, which inevitably causes depth bleeding across object boundaries. To address this limitation, we design the Lambertian EPI Branch. Its core objective is to extract initial depth from EPI line structures and subsequently refine it using a continuous depth prior, specifically the plane regularized sampling operator (PRSO), to guarantee piecewise smoothness while preserving sharp depth discontinuities.

Specifically, the module takes the masked Lambertian light field data $L_L(x,y,u,v)$ and the pixel-level mask $M(x,y)$ as inputs. The spatial dimensions are $H \times W$ and the angular dimensions are $U \times V$. We first extract EPI features to regress the initial disparity map. Let $\mathcal{F}_{3D}$ denote a 3D convolutional feature extractor that captures the spatio-angular correlations along the epipolar lines. The initial disparity $D_{init}(x,y)$ is obtained through a slope regression head $\mathcal{R}_{slope}$, formulated as:
\begin{equation} \label{eq:ch3_8}
D_{init}(x,y) = \mathcal{R}_{slope}(\mathcal{F}_{3D}(L_L(x,y,u,v)))
\end{equation}
where $D_{init}(x,y)$ represents the preliminary depth estimate derived purely from local EPI gradients. While accurate in highly textured areas, $D_{init}(x,y)$ suffers from severe degradation in textureless regions due to the lack of distinguishable angular variations.

To recover reliable depth in these under-constrained regions, we introduce a continuous depth prior optimization mechanism driven by the PRSO. Instead of enforcing global smoothness, we model the local depth surface as a continuous 3D plane. For a given spatial pixel $(x,y)$, we define a local neighborhood $\mathcal{N}_r(x,y)$ with radius $r$. Within this neighborhood, the local depth is approximated by a plane parameterized by $\mathbf{p}_{x,y} = [a, b, c]^T$, such that the planar depth $\hat{D}(x',y'; \mathbf{p}_{x,y})$ for any pixel $(x',y') \in \mathcal{N}_r(x,y)$ is given by:
\begin{equation} \label{eq:ch3_9}
\hat{D}(x',y'; \mathbf{p}_{x,y}) = a x' + b y' + c
\end{equation}
where $a$, $b$, and $c$ represent the plane normal components and the offset, respectively. The plane parameters $\mathbf{p}_{x,y}$ are estimated via weighted least squares fitting using the reliable initial disparities in the local neighborhood.

We then formulate the depth refinement as an energy minimization problem to compute the final Lambertian depth map $D_L(x,y)$. The objective function integrates a data fidelity term and a local plane regularization term, explicitly guided by the mask $M(x,y)$:
\begin{equation} \label{eq:ch3_10}
E(D_L) = \sum_{(x,y)} M(x,y) \left[ \lambda_1 \rho(D_L(x,y) - D_{init}(x,y)) + \lambda_2 \sum_{(x',y') \in \mathcal{N}_r(x,y)} \omega(x',y') \rho(D_L(x',y') - \hat{D}(x',y'; \mathbf{p}_{x,y})) \right]
\end{equation}
where $\lambda_1$ and $\lambda_2$ are balancing weights controlling the trade-off between data fidelity and spatial regularization. The function $\rho(\cdot)$ denotes the Charbonnier penalty function, which provides robustness against outliers in the initial regression. The term $\omega(x',y')$ represents a bilateral weight computed from spatial distance and photometric similarity, ensuring that the plane regularization respects local image structures.

Different from previous works that apply uniform global regularization causing depth artifacts at edges, our module explicitly incorporates the pixel-level mask $M(x,y)$ to restrict the optimization domain strictly to Lambertian regions. This prevents the anomalous gradients from non-Lambertian reflections from corrupting the local plane fitting. Furthermore, by leveraging the PRSO-based local plane prior, the module effectively interpolates depth in textureless areas while maintaining piecewise smoothness. The optimization yields the refined Lambertian depth map $D_L(x,y)$ with dimensions $H \times W$. Finally, $D_L(x,y)$ is forwarded to the subsequent Dual-Mask Fusion module, where it will be integrated with the non-Lambertian depth estimates to produce the unified final depth map.

\subsection{Non-Lambertian Geometric Branch}

Following the extraction of the Lambertian depth map $D_L$ in the preceding module, we route the remaining light field data to the Non-Lambertian Geometric Branch. Unlike Lambertian surfaces that exhibit clear linear structures in epipolar plane image (EPI)s, Non-Lambertian surfaces (e.g., specular highlights and reflections) manifest complex X-shaped patterns and bimodal distributions along the angular dimensions. Previous light field depth estimation methods typically treat these Non-Lambertian regions as noise or outliers, applying global smoothness priors that inevitably cause severe depth blurring and structural degradation. To address this limitation, we design the Non-Lambertian Geometric Branch to explicitly exploit these unique geometric characteristics rather than suppress them. The core objective of this module is to accurately estimate the depth map $D_{NL}$ for Non-Lambertian regions by leveraging a cross-guided architecture and bimodal distribution feature matching, thereby providing reliable depth cues for the subsequent Dual-Mask Fusion module.

Specifically, the module takes the Non-Lambertian light field data $L_{NL}(x,y,u,v)$, the geometric features, and the pixel-level mask $M(x,y)$ as inputs, where $(x,y)$ denotes the spatial coordinates and $(u,v)$ represents the angular coordinates. The spatial and angular dimensions are $H \times W$ and $U \times V$, respectively. Since Non-Lambertian regions, particularly specular highlights, suffer from severe data scarcity in existing annotated datasets, direct training often leads to overfitting. To mitigate this issue, we introduce a tone-mapping data augmentation strategy to simulate diverse illumination conditions and enrich the training distribution. We apply a randomized gamma correction to the input light field, formulated as:
\begin{equation} \label{eq:ch3_11}
L_{NL}^{\gamma}(x,y,u,v) = \left( \frac{L_{NL}(x,y,u,v)}{L_{max}} \right)^{\gamma} \cdot L_{max}
\end{equation}
where $L_{max}$ is the maximum intensity value of the light field, and $\gamma$ is a random sampling parameter drawn from a uniform distribution $\mathcal{U}(\gamma_{min}, \gamma_{max})$. This augmentation effectively expands the feature space of Non-Lambertian reflections without altering the underlying geometric structures.

To comprehensively capture both the local X-shaped patterns and the global bimodal distributions, we propose a dual-branch cross-guided architecture. Different from previous works that rely on a single feature extractor, our module explicitly separates the feature extraction into a local pattern branch and a global distribution branch. Let $\mathcal{E}_{loc}$ and $\mathcal{E}_{glob}$ denote the local and global feature extractors, respectively. We first extract the initial feature maps $\mathbf{F}_{loc} = \mathcal{E}_{loc}(L_{NL}^{\gamma})$ and $\mathbf{F}_{glob} = \mathcal{E}_{glob}(L_{NL}^{\gamma})$, where $\mathbf{F}_{loc}, \mathbf{F}_{glob} \in \mathbb{R}^{H \times W \times C}$ and $C$ is the channel dimension. To facilitate information exchange between the two branches, we employ a cross-guided attention mechanism. The refined local features $\mathbf{F}_{loc}'$ are updated by attending to the global features, and vice versa:
\begin{equation} \label{eq:ch3_12}
\mathbf{F}_{loc}' = \mathbf{F}_{loc} + \text{Softmax}\left(\frac{\mathbf{Q}_{loc}\mathbf{K}_{glob}^T}{\sqrt{d_k}}\right)\mathbf{V}_{glob}
\end{equation}
\begin{equation} \label{eq:ch3_13}
\mathbf{F}_{glob}' = \mathbf{F}_{glob} + \text{Softmax}\left(\frac{\mathbf{Q}_{glob}\mathbf{K}_{loc}^T}{\sqrt{d_k}}\right)\mathbf{V}_{loc}
\end{equation}
where $\mathbf{Q}$, $\mathbf{K}$, and $\mathbf{V}$ represent the query, key, and value projections derived from the respective feature maps, and $d_k$ is the scaling factor. This cross-guided operation ensures that the local structural details are regularized by global context, while the global distribution is refined by local geometric cues.

Following the feature refinement, we perform bimodal distribution geometric feature matching to regress the final depth. In Non-Lambertian regions, the angular signal profile of a pixel typically exhibits a bimodal distribution, consisting of a diffuse reflection peak and a specular reflection peak. We define the angular profile $\mathbf{p}_{x,y} \in \mathbb{R}^{U \times V}$ for each spatial location. To measure the matching cost across different disparity hypotheses $d$, we construct a cost volume $\mathbf{V}_{cost} \in \mathbb{R}^{H \times W \times D_{max}}$, where $D_{max}$ is the maximum disparity. The cost volume is computed by evaluating the difference between the shifted angular profiles:
\begin{equation} \label{eq:ch3_14}
\mathbf{V}_{cost}(x,y,d) = \sum_{u,v} \left\| \mathbf{p}_{x,y}(u,v) - \mathbf{p}_{x-u \cdot d, y-v \cdot d}(u,v) \right\|_2^2
\end{equation}
To explicitly leverage the bimodal property, we calculate a bimodality weight map $\mathbf{W}_{bi}(x,y)$ based on the bimodality coefficient of the angular profile $\mathbf{p}_{x,y}$. This weight map assigns higher confidence to pixels with distinct bimodal distributions. The final Non-Lambertian depth map $D_{NL}(x,y)$ is then obtained by applying a soft argmin operation over the weighted cost volume:
\begin{equation} \label{eq:ch3_15}
D_{NL}(x,y) = \sum_{d=0}^{D_{max}} d \cdot \frac{\exp(-\mathbf{V}_{cost}(x,y,d) \cdot \mathbf{W}_{bi}(x,y))}{\sum_{k=0}^{D_{max}} \exp(-\mathbf{V}_{cost}(x,y,k) \cdot \mathbf{W}_{bi}(x,y))}
\end{equation}
This formulation guarantees that the depth estimation is heavily guided by the reliable bimodal geometric features, effectively suppressing the interference from ambiguous specular highlights.

The output of this module is the Non-Lambertian depth map $D_{NL} \in \mathbb{R}^{H \times W}$, which preserves sharp depth discontinuities and accurate structural details in reflective regions. Together with the Lambertian depth map $D_L$ and the pixel-level mask $M(x,y)$ from the preceding modules, $D_{NL}$ is subsequently fed into the Dual-Mask Fusion module. By explicitly modeling the physical and geometric properties of Non-Lambertian surfaces, our proposed branch effectively overcomes the limitations of traditional EPI-based assumptions, providing a robust and unified depth estimation framework.

\subsection{Dual-Mask Fusion}

Following the extraction of the specialized depth maps from the Lambertian EPI Branch and the Non-Lambertian Geometric Branch, we obtain $D_L$ and $D_{NL}$, respectively. While these two branches accurately capture the depth structures within their specific physical domains, a unified depth representation is required for downstream applications. Existing light field depth estimation methods typically merge multi-domain results using post-processing global smoothness priors or heuristic thresholding. Such strategies inevitably cause depth bleeding across domain boundaries and structural degradation, primarily because they ignore the distinct physical properties of the surfaces. To address this limitation, we design the Dual-Mask Fusion module to effectively integrate the dual-branch depth estimations guided by the physically derived pixel-level mask.

The core principle of this module involves utilizing the pixel-level geometric mask to perform a domain-aware weighted fusion, thereby eliminating the errors caused by the failure of the single EPI assumption in Non-Lambertian regions. Specifically, the module takes the Lambertian depth map $D_L \in \mathbb{R}^{H \times W}$, the Non-Lambertian depth map $D_{NL} \in \mathbb{R}^{H \times W}$, and the pixel-level mask $M \in [0, 1]^{H \times W}$ as inputs. Here, $H \times W$ denotes the spatial resolution of the light field, and $M(x,y)$ represents the probability or binary indicator of the pixel at spatial coordinates $(x,y)$ belonging to the Lambertian domain, which is generated by the preceding Geometric Feature \& Mask Generation module.

To formulate the fusion process, we aim to minimize the domain-specific estimation error at each spatial location. Let $E_L(x,y)$ and $E_{NL}(x,y)$ denote the hypothetical depth errors of the Lambertian and Non-Lambertian branches, respectively. Based on the physical decomposition in the Three-Layer Angular Signal Decomposition module, the Lambertian branch yields minimal error in diffuse regions where $M(x,y) \to 1$, while the Non-Lambertian branch is optimal in specular regions where $M(x,y) \to 0$. Therefore, we construct a pixel-wise objective to minimize the expected error, leading to a convex combination of the two branch outputs. We compute the final unified depth map $D_{final} \in \mathbb{R}^{H \times W}$ through the following mathematical formulation:
\begin{equation} \label{eq:ch3_16}
D_{final}(x,y) = M(x,y) \odot D_L(x,y) + (1 - M(x,y)) \odot D_{NL}(x,y)
\end{equation}
where $\odot$ denotes the Hadamard product (element-wise multiplication). The term $M(x,y)$ acts as a spatial attention weight that selectively preserves the reliable depth cues from the Lambertian EPI Branch for diffuse surfaces. Conversely, the complementary weight $(1 - M(x,y))$ activates the depth estimations from the Non-Lambertian Geometric Branch for specular and reflective regions. 

Different from previous works that apply uniform regularization across the entire image, our module explicitly respects the physical domain boundaries defined by $M(x,y)$. This design prevents the inappropriate smoothing of Non-Lambertian depth discontinuities by Lambertian priors. Furthermore, by directly routing the depth values based on the geometric mask, we avoid the computational overhead and potential sub-optimal convergence associated with complex learnable fusion networks. 

Ultimately, the Dual-Mask Fusion module outputs the unified depth map $D_{final}$, which retains the piecewise smoothness in Lambertian regions while preserving the sharp geometric structures in Non-Lambertian areas. This unified representation not only resolves the domain conflict inherent in mixed-reflectance scenes but also provides a robust and accurate depth estimation for the entire light field.

\subsection{Training Objective}

\subsubsection{Training Objective}

Training a unified depth estimation model across Lambertian, mixed, and non-Lambertian domains presents a severe data imbalance challenge. The non-Lambertian dataset contains only four scenes, rendering the model highly susceptible to overfitting on these rare samples while underfitting the complex texture variations. Furthermore, the Epipolar Plane Image (EPI) assumption inherently degrades on specular and scattering surfaces. Therefore, a standard mean absolute error (MAE) objective proves insufficient. We propose a composite training objective that jointly optimizes depth regression, material-aware mask generation, and physical consistency to explicitly guide the GeometricDualMask network.

\paragraph{Overall Objective}
The total loss function $\mathcal{L}_{total}$ is formulated as a weighted sum of three principal components:
\begin{equation} \label{eq:ch3_17}
\mathcal{L}_{total} = \lambda_{depth} \mathcal{L}_{depth} + \lambda_{mask} \mathcal{L}_{mask} + \lambda_{reg} \mathcal{L}_{reg}
\end{equation}
where $\mathcal{L}_{depth}$ denotes the depth regression loss, $\mathcal{L}_{mask}$ represents the dual-mask classification loss, and $\mathcal{L}_{reg}$ is the physical consistency regularization term. The scalars $\lambda_{depth}$, $\lambda_{mask}$, and $\lambda_{reg}$ balance the contributions of each objective during the joint optimization.

\paragraph{Depth Estimation Loss}
\textbf{Motivation:} Standard $L_1$ or $L_2$ losses treat all pixels equally, which causes the network to ignore hard examples, such as depth boundaries and non-Lambertian highlights where EPI slopes are severely distorted. 

\textbf{Design:} We employ a Focal Smooth $L_1$ loss to down-weight well-predicted Lambertian regions and focus optimization on the challenging non-Lambertian pixels. Let $D$ and $\hat{D}$ be the ground-truth and predicted depth maps, respectively. The Smooth $L_1$ loss for a pixel $i$ is defined as:
\begin{equation} \label{eq:ch3_18}
\ell_{smooth}(d_i, \hat{d}_i) = \begin{cases} 0.5 (d_i - \hat{d}_i)^2 & \text{if } |d_i - \hat{d}_i| < 1 \\ |d_i - \hat{d}_i| - 0.5 & \text{otherwise} \end{cases}
\end{equation}
To incorporate the focal mechanism, we modulate this loss with a focusing parameter $\gamma$:
\begin{equation} \label{eq:ch3_19}
\mathcal{L}_{depth} = \frac{1}{N} \sum_{i=1}^{N} (1 - p_i)^\gamma \ell_{smooth}(d_i, \hat{d}_i)
\end{equation}
where $N$ is the total number of pixels, and $p_i = \exp(-\ell_{smooth}(d_i, \hat{d}_i))$ acts as a pseudo-probability of correct prediction. The parameter $\gamma > 0$ controls the down-weighting rate. 

\textbf{Expected Effect:} This design explicitly forces the GeometricDualMask architecture to allocate more representational capacity to resolving the ambiguous EPI structures in non-Lambertian regions, thereby improving the overall MAE on mixed and non-Lambertian domains.

\paragraph{Dual-Mask Classification Loss}
\textbf{Motivation:} The GeometricDualMask module relies on accurate pixel-level material masks to route features into domain-specific branches. However, non-Lambertian pixels (e.g., specular highlights) occupy a minor fraction of the image, leading to severe class imbalance.

\textbf{Design:} Different from previous works that utilize standard cross-entropy for mask prediction, our module explicitly combines Focal Loss and Dice Loss to address this extreme pixel-level imbalance. Let $M$ and $\hat{M}$ be the ground-truth and predicted binary masks (1 for non-Lambertian, 0 for Lambertian). The Dice loss $\mathcal{L}_{dice}$ maximizes the overlap between prediction and ground truth:
\begin{equation} \label{eq:ch3_20}
\mathcal{L}_{dice} = 1 - \frac{2 \sum_{i} m_i \hat{m}_i + \epsilon}{\sum_{i} m_i + \sum_{i} \hat{m}_i + \epsilon}
\end{equation}
where $\epsilon$ is a small constant to prevent division by zero. The mask loss is then formulated as:
\begin{equation} \label{eq:ch3_21}
\mathcal{L}_{mask} = \alpha \mathcal{L}_{focal\_mask} + (1 - \alpha) \mathcal{L}_{dice}
\end{equation}
where $\mathcal{L}_{focal\_mask}$ is the standard binary focal loss applied to the mask logits, and $\alpha \in [0, 1]$ balances the two terms. 

\textbf{Expected Effect:} This composite loss ensures high recall for rare non-Lambertian pixels while maintaining precise boundary delineation, preventing the mask generator from collapsing into a trivial all-Lambertian prediction.

\paragraph{Physical Consistency Regularization}
\textbf{Motivation:} In non-Lambertian regions, the linear EPI slope is corrupted by the residual material component. Relying solely on data-driven depth supervision is inadequate due to the scarcity of non-Lambertian training scenes. We must enforce physical constraints derived from our three-layer angular signal decomposition.

\textbf{Design and Derivation:} Recall that the angular signal at pixel $i$ is decomposed as $S_i = S_i^{DC} + S_i^{linear} + S_i^{residual}$, representing ambient illumination, disparity, and material reflectance, respectively. The linear component $S_i^{linear}$ encodes the epipolar slope, which is directly proportional to the spatial depth gradient $\nabla \hat{d}_i$. To isolate the geometric structure from the material reflectance (encoded in $S_i^{residual}$), we constrain the variation of $S_i^{linear}$ in Lambertian regions where $S_i^{residual} \approx 0$. Mathematically, minimizing the variation of the linear component translates to penalizing the depth gradients weighted by the inverse of the image edges and the Lambertian mask. This yields the edge-aware regularization formulation:
\begin{equation} \label{eq:ch3_22}
\mathcal{L}_{reg} = \frac{1}{N} \sum_{i} \left( |\nabla_x \hat{d}_i| e^{-\beta |\nabla_x I_i|} (1 - \hat{m}_i) + |\nabla_y \hat{d}_i| e^{-\beta |\nabla_y I_i|} (1 - \hat{m}_i) \right)
\end{equation}
where $\nabla_x$ and $\nabla_y$ denote spatial gradients, $I$ is the central view image, and $\beta$ is a sensitivity parameter. The term $(1 - \hat{m}_i)$ explicitly masks out the non-Lambertian regions from the smoothness penalty.

\textbf{Expected Effect:} By acknowledging that depth estimation in non-Lambertian regions relies on the specialized branch rather than local spatial smoothness, this physical constraint effectively decouples geometry from reflectance, stabilizing the training process under data-scarce conditions.

\paragraph{Weighting Strategy and Domain-Balanced Optimization}
\textbf{Motivation:} The extreme scarcity of non-Lambertian data necessitates careful loss weighting and sampling strategies to prevent catastrophic forgetting of the abundant Lambertian and mixed domains.

\textbf{Design:} We set the loss weights empirically based on validation performance: $\lambda_{depth} = 1.0$, $\lambda_{mask} = 0.5$, and $\lambda_{reg} = 0.1$. The lower weight for $\mathcal{L}_{reg}$ prevents the smoothness constraint from over-smoothing fine geometric details. For $\mathcal{L}_{mask}$, we set $\alpha = 0.5$ and $\gamma = 2.0$ for the focal components. Furthermore, to complement the loss function, we employ a Domain-balanced sampling strategy during training. A WeightedRandomSampler assigns higher sampling probabilities to the non-Lambertian and mixed domains, ensuring that each mini-batch contains a balanced distribution of physical reflections. 

\textbf{Expected Effect:} This joint optimization of the composite loss and balanced sampling effectively mitigates the domain shift, alleviates overfitting on the minority non-Lambertian scenes, and enables the unified model to achieve robust depth estimation across diverse material properties.

\section{Experiments}

\subsection{Datasets}

To comprehensively evaluate the proposed Unified Dual-Mask Physical Model across diverse reflectance properties and scene complexities, we construct a multi-domain light field (LF) benchmark. The selected datasets encompass a wide spectrum of surface characteristics, ranging from ideal Lambertian to highly challenging non-Lambertian (e.g., specular, scattering, and transparent) materials. Specifically, our experiments are conducted on three primary datasets: the HCI New dataset \cite{visionappl}, the UrbanLF-Syn dataset \cite{nonlamb}, and a curated Non-Lambertian dataset.

\subsubsection{Dataset Descriptions}

\textbf{HCI New Dataset.} The HCI New dataset \cite{visionappl} is a widely recognized synthetic benchmark primarily consisting of Lambertian scenes. It includes multiple complex scenes (e.g., \textit{boxes}, \textit{cotton}, \textit{dino}, and \textit{stratified}) with rich textures and intricate geometric structures. The dataset provides an angular resolution of $9 \times 9$ (81 views) and a spatial resolution of $512 \times 512$ pixels. The ground truth depth maps are perfectly generated by the rendering engine, providing pixel-accurate supervision. We utilize the official training and validation splits to assess the model's foundational depth estimation capability under ideal Lambertian assumptions.

\textbf{UrbanLF-Syn Dataset.} To evaluate the model's performance in mixed and complex urban environments, we employ the UrbanLF-Syn dataset \cite{nonlamb}. This dataset comprises 170 synthetic scenes featuring diverse urban structures, varying illumination conditions, and mixed surface reflectance (i.e., a combination of Lambertian and mild non-Lambertian regions). The angular and spatial resolutions are standardized to $9 \times 9$ and $512 \times 512$, respectively. The large number of scenes provides robust training data for the mixed/urban domain, and the depth ground truth is derived from the synthetic rendering pipeline.

\textbf{Non-Lambertian Dataset.} Evaluating depth estimation on non-Lambertian surfaces remains a formidable challenge due to the violation of the photo-consistency assumption in Epipolar Plane Images (EPIs). For this domain, we utilize a specialized Non-Lambertian dataset featuring severe specular highlights, scattering, and transparent/mirror (ToM) surfaces. The scenes are synthesized utilizing measured Bidirectional Reflectance Distribution Functions (BRDFs) from the MERL BRDF database \cite{epistru} to ensure physically accurate complex reflectance. Crucially, this dataset suffers from severe data scarcity, containing only 4 training scenes. The angular and spatial resolutions are $9 \times 9$ and $512 \times 512$. The extreme scarcity of non-Lambertian LF data constitutes a significant domain imbalance, serving as a rigorous testbed for the generalization and physical modeling capabilities of our unified framework.

\subsubsection{Dataset Selection Rationale and Challenges}

The selection of these three datasets is deliberately designed to cover the full trajectory of light field depth estimation challenges. The \textit{HCI New} dataset serves as the baseline for evaluating the upper bound of performance under standard Lambertian conditions. The \textit{UrbanLF-Syn} dataset introduces real-world complexity and mixed materials, testing the model's robustness to moderate domain shifts and occlusions. 

The most critical challenge lies in the \textit{Non-Lambertian} dataset. In non-Lambertian regions, the fundamental EPI linearity assumption breaks down, rendering traditional and even many deep learning-based EPI architectures ineffective. Furthermore, the extreme data scarcity (only 4 scenes) creates a severe long-tail distribution across domains. This intentional inclusion allows us to rigorously validate whether the proposed Dual-Mask Physical Model and Angular FFT Material Classification can effectively decouple domain-specific physical features and prevent the model from overfitting to the dominant Lambertian/Mixed domains.

\subsubsection{Data Preprocessing and Sampling Strategy}

To ensure consistency across the multi-domain benchmark, all light field images are preprocessed to a uniform spatial resolution of $512 \times 512$ and an angular resolution of $9 \times 9$. Depth maps are normalized to a unified range to facilitate stable joint optimization of the depth regression and domain classification objectives.

To mitigate the severe domain imbalance---particularly the drastic disparity between the 170 UrbanLF scenes and the 4 Non-Lambertian scenes---we implement a \textbf{Domain-Balanced Sampling} strategy during training. Specifically, we employ a \texttt{WeightedRandomSampler} that dynamically assigns higher sampling probabilities to the minority domains (Non-Lambertian) and lower probabilities to the majority domains (UrbanLF-Syn). This ensures that each training batch contains a balanced representation of all physical domains, preventing the network from being biased toward the Lambertian/Mixed priors and encouraging the dual-mask modules to learn domain-invariant and domain-specific physical features effectively. The datasets are strictly partitioned into training and validation sets, with no overlap in scene identities to ensure a fair evaluation of generalization.

\begin{table*}[!t]
\caption{Summary of Datasets}
\label{tab:datasets}
\centering
\begin{tabular}{llcccll}
\toprule
\textbf{Dataset} & \textbf{Scene Type} & \textbf{No. of Scenes} & \textbf{Angular Res.} & \textbf{Spatial Res.} & \textbf{Depth Source} & \textbf{Key Characteristics} \\
\midrule
\textbf{HCI New} & Lambertian & 4 train + 4 val & $9 \times 9$ & $512 \times 512$ & Synthetic GT & Ideal photo-consistency, rich textures, standard benchmark. \\
\textbf{UrbanLF-Syn} & Mixed / Urban & 170 & $9 \times 9$ & $512 \times 512$ & Synthetic GT & Complex urban structures, mixed reflectance, large scale. \\
\textbf{Non-Lambertian} & Non-Lambertian & 4 & $9 \times 9$ & $512 \times 512$ & Synthetic GT (MERL BRDF) & Severe specular/scattering, EPI assumption breakdown, extreme data scarcity. \\
\bottomrule
\end{tabular}
\vspace{2mm}

\footnotesize{\textit{*Note: The "No. of Scenes" for HCI New refers to the specific subsets (e.g., training and stratified) used in our experiments. All datasets are strictly split into training and validation sets with no overlap.}}
\end{table*}

\subsection{Implementation Details}

\subsubsection{Hardware and Software Environment}

All experiments were implemented using the PyTorch 2.0 framework and trained on a high-performance computing cluster equipped with four NVIDIA RTX A6000 GPUs (48GB VRAM each). The software environment was built on CUDA 11.8 and cuDNN 8.6 to accelerate tensor operations. To maximize hardware utilization and synchronize gradients across multiple devices, Distributed Data Parallel (DDP) training was employed. 

\subsubsection{Training Hyperparameters and Optimization}

The proposed Unified Dual-Mask Physical Model is built upon the EPINet4Dir V3 backbone, which extracts 4-directional Epipolar Plane Images (EPIs). We employ the AdamW optimizer with a weight decay of $1 \times 10^{-4}$ to mitigate overfitting, particularly in the data-scarce Non-Lambertian domain. The initial learning rate is set to $1 \times 10^{-4}$, which is dynamically adjusted using a Cosine Annealing learning rate scheduler with a minimum learning rate of $1 \times 10^{-6}$ and a warm-up phase of 5 epochs. 

Due to the substantial memory footprint of 4D light field data (e.g., $9 \times 9$ angular views at $512 \times 512$ spatial resolution), the per-GPU batch size is restricted to 4. To maintain a sufficient effective batch size for stable gradient estimation, we utilize gradient accumulation with an accumulation step of 4, yielding an effective batch size of 64 across the four GPUs. The model is trained for a maximum of 300 epochs. Early stopping with a patience of 30 epochs is applied based on the validation Mean Absolute Error (MAE) to prevent overfitting. The key training hyperparameters are summarized in Table \ref{tab:hyperparameters}.

\begin{table}[!t]
\caption{Summary of Key Training Hyperparameters}
\label{tab:hyperparameters}
\centering
\begin{tabular}{llll}
\toprule
\textbf{Hyperparameter} & \textbf{Configuration} & \textbf{Hyperparameter} & \textbf{Configuration} \\
\midrule
\textbf{Optimizer} & AdamW & \textbf{Base Model} & EPINet4Dir V3 \\
\textbf{Initial Learning Rate} & $1 \times 10^{-4}$ & \textbf{Weight Decay} & $1 \times 10^{-4}$ \\
\textbf{Min Learning Rate} & $1 \times 10^{-6}$ & \textbf{Warm-up Epochs} & 5 \\
\textbf{Batch Size (per GPU)} & 4 & \textbf{Gradient Accum.} & 4 \\
\textbf{Effective Batch Size} & 64 & \textbf{Max Epochs} & 300 \\
\textbf{Early Stopping} & Patience 30 & \textbf{Scheduler} & Cosine Annealing \\
\bottomrule
\end{tabular}
\end{table}

Building upon these comprehensive empirical validations, we now distill our core contributions and outline the broader impact of this work.

\section{Conclusion}

In this paper, we have proposed a unified dual-mask physical model to address the persistent challenge of depth estimation for non-Lambertian surfaces in light field imaging, where the fundamental Epipolar Plane Image (EPI) linearity assumption is severely violated. By explicitly modeling the physical degradation of angular signals, our framework decouples material-specific reflections from geometric structures, enabling robust depth inference in complex mixed-reflectance environments. 

The main contributions of this work are summarized as follows:
\begin{itemize}
    \item \textbf{Geometry-based Material Classification via Angular Signal Decomposition:} We introduced a three-layer decomposition mechanism that overcomes frequency aliasing caused by discrete low-resolution angular sampling, providing a robust and efficient material identification scheme that lays the foundation for dual-branch depth estimation.
    \item \textbf{Unified Dual-Mask Depth Estimation Architecture:} We developed a dual-branch framework that effectively mitigates the depth collapse typically encountered by single-EPI-based methods when facing specular or translucent surfaces due to physical assumption failures.
    \item \textbf{Mask Routing for Mixed-Scene Generalization:} By integrating a dynamic mask routing mechanism within a unified architecture, the proposed model achieves high-precision depth estimation across both mixed and holistic scenes, significantly improving the generalization capability and practical applicability of the system.
\end{itemize}

Although the proposed unified architecture achieves robust performance in mixed scenarios, it remains fundamentally constrained by the extreme scarcity of non-Lambertian training data and the inherent blurring of EPI slopes in complex textured regions. To transcend these intrinsic bottlenecks of epipolar geometry, future research will shift towards non-EPI paradigms, such as Neural Radiance Fields (NeRF), leveraging substantially expanded multi-reflectance datasets to achieve comprehensive depth estimation across all material domains.

The explicit three-layer angular signal decomposition offers a distinct advantage over purely data-driven, black-box convolutional neural networks (CNNs) by enforcing physical priors directly within the signal processing pipeline. By decoupling the angular signal $I(u,v)$ into the DC component (ambient illumination), the linear component (epipolar geometry), and the residual $\epsilon(u,v)$ (material reflectance), our model prevents the entanglement of high-frequency specular highlights with low-frequency disparity cues. In end-to-end architectures, the network often struggles to disentangle these overlapping frequency bands, leading to depth collapse in non-Lambertian regions where the Bidirectional Reflectance Distribution Function (BRDF) varies rapidly. Our physics-driven decoupling ensures that the residual $\epsilon(u,v)$ safely absorbs material-induced radiometric variations, guaranteeing that the geometric disparity estimation remains strictly invariant to surface reflectance properties, thereby enhancing both algorithmic robustness and physical interpretability.

A critical design choice in our framework is the deliberate departure from frequency-domain analyses, such as the 2D Discrete Fourier Transform (DFT), in favor of spatial-domain geometric features. From a signal sampling perspective, commercial plenoptic cameras typically provide low angular resolutions (e.g., $9 \times 9$ views). According to the Nyquist-Shannon sampling theorem, such sparse angular sampling inevitably induces severe spectral aliasing and leakage in the frequency domain, rendering DFT-based material classification fundamentally unreliable. By extracting spatial geometric features---such as angular gradients, correlation coefficients, and coefficients of variation---we effectively bypass the frequency aliasing bottleneck. These geometric descriptors act as robust spatial-domain statistics that preserve structural integrity without requiring dense angular sampling. The high Cohen's $d$ effect size of the angular gradient feature empirically validates that spatial-domain metrics remain highly discriminative for material identification even under severely discrete angular sampling conditions.

Furthermore, the proposed unified dual-mask routing mechanism addresses the inherent gradient conflict encountered when optimizing a single monolithic network for mixed-reflectance scenes. In conventional single-branch models, the loss gradients derived from Lambertian and non-Lambertian regions often point in opposing directions in the parameter space, causing sub-optimal convergence and feature suppression. By dynamically routing features to specialized sub-networks based on the material masks, we effectively isolate the optimization landscapes. Computationally, utilizing a lightweight Random Forest (RF) classifier for the initial mask generation is significantly more efficient than deploying heavy, pixel-wise semantic segmentation networks. However, this cascaded architecture introduces a potential vulnerability: misclassifications in the RF stage could propagate as erroneous hard masks. To mitigate this, the domain-balanced sampling strategy and the inherent tolerance of the multi-directional EPI features ensure that minor mask boundary inaccuracies do not catastrophically degrade the final depth maps. Future iterations could explore differentiable soft-mask routing to enable unified end-to-end gradient flow.

Despite the demonstrated efficacy, the current physical model operates under specific boundary conditions. The linear approximation of the epipolar component assumes relatively small baselines and opaque or semi-translucent surfaces. For highly refractive transparent objects (e.g., thick glass or fluids), light rays undergo complex bending, fundamentally violating the straight-line EPI assumption and producing highly non-linear, fragmented trajectories. Addressing such extreme optical phenomena would require integrating explicit ray-tracing priors or transitioning to volumetric rendering paradigms, such as Neural Radiance Fields (NeRF) or 3D Gaussian Splatting, specifically adapted for multi-view light field inputs. Additionally, extending the current discrete angular decomposition to continuous light field representations could unlock higher-fidelity depth estimation for next-generation plenoptic sensors, while incorporating temporal consistency constraints would be essential for deploying this framework in dynamic, real-world robotic navigation scenarios.

While the proposed unified dual-mask framework demonstrates robust performance in decoupling angular signals, several physical and computational limitations remain, outlining critical directions for future research.

First, the three-layer angular signal decomposition relies on a first-order local linearity assumption, modeled as $I(u,v) = I_{dc} + a \cdot u + b \cdot v + \epsilon(u,v)$. While this effectively isolates material residuals $\epsilon(u,v)$ for standard specular and scattering surfaces, it exhibits limitations under complex global illumination, such as inter-reflections or caustics. In such regimes, the residual term becomes non-linearly coupled with geometric disparity, potentially degrading the Random Forest classifier's precision and inducing mask routing artifacts at high-contrast material boundaries. Future work could integrate higher-order bidirectional reflectance distribution function (BRDF) models or hybridize our explicit physical decomposition with implicit neural rendering to better disentangle multi-bounce light transport.

Second, although our geometry-based feature extraction circumvents frequency aliasing at low angular resolutions (e.g., $9 \times 9$), the discrete spatial sampling of the Epipolar Plane Image (EPI) restricts the recovery of sub-pixel or extremely fine structures (e.g., hair, thin wires). At these scales, EPI linearity inherently degrades due to spatial aliasing, leading to potential depth discontinuities. To mitigate this, future iterations could incorporate spatial super-resolution priors or exploit sub-pixel phase correlation to enhance EPI structural integrity prior to dual-mask routing.

Finally, the current signal processing pipeline is strictly formulated for static light fields. Extending the architecture to spatiotemporal light fields introduces temporal motion blur, which acts as an additional low-pass filter on angular signals, further obscuring material-specific high-frequency residuals. Future research will investigate integrating temporal consistency constraints into the mask routing mechanism. Furthermore, from a hardware-algorithm co-design perspective, leveraging polarized light field sensors or event-based vision could provide direct physical measurements of surface normals and material properties. This would bypass the computational overhead of algorithmic signal decoupling, enabling real-time, high-fidelity depth estimation in dynamic, mixed-reflectance environments.

\bibliographystyle{IEEEtran}
\bibliography{references}

\end{document}
